%% Unit 4 -- Verify without reading every line.  (Ch6 Implementation, Tests, Review.)
%% Restyled to the sparse, one-idea-per-slide house style of units 1-3: one
%% thesis-shaped title, prose-first bodies, EXACTLY one red per rendered frame
%% (the cmdtable frame keeps its \gcmd slash-commands as identity; a red
%% transcript OUTCOME line counts as that frame's one red). Scientific claims are
%% cited via \slidesources; the closing frame prints the unit's papers. Realized
%% visual-lean: the risk graph and the deep-vs-shallow contrast carry their own
%% frames. No preamble here; \input by the master (disciplined_vibe.tex).
%%
%% ORDER follows the build half of the software-engineering process:
%%   A. Handoff      -- the day shift is done; queue the build.
%%   B. Build skills -- TDD, code-review, diagnosing-bugs,
%%                      improve-codebase-architecture, write-a-skill.
%%   C. Breakout     -- vibe the banana again, now tested.

\sectionsubtitle{aber mit höchster Disziplin.}
% ======================================================================
\section[Building]{Software Engineering Knows How already}
% ======================================================================


% ======================================================================
% A. HANDOFF -- the day shift is done; queue the build.
% ======================================================================

% === Now let the agent build one ===
% A transition slide: one idea, no margin column. The takeaway is a full-width
% centred display (it \vfill-s itself), so we let it carry the frame rather than
% cramming it into a tdmain beside an empty tdmargin.
\begin{frame}{And so we build}
  \vskip0.8em
  That was the \alertbf{day shift}: plan, align, and slice by hand,
  human-in-the-loop. You now hold a backlog of trustworthy slices molded into issues.\\
  \vskip0.8em
  You just finished a product owner workflow, ready to hand-off to the Devs.
  Remember the silverbullet.
\end{frame}


\iffalse
% ======================================================================
% B. BUILD SKILLS
% ======================================================================

% === the build toolkit: the read-only catalogue this unit walks, one frame ===
% Mirrors the plan toolkit map in Unit 3: same four-column read-only layout,
% now for the build half. Only the build skills live here; the plan arc
% (askuserquestion, domain-modeling, to-prd/to-issues, prototype) belongs to
% Unit 3, so keep it off this frame so it maps only what Unit 4 covers.
\begin{frame}{The build toolkit, read only}
  \vskip0.3em
  \begin{columns}[T,onlytextwidth]
    \begin{column}{0.23\textwidth}
      \begin{tdblock}
        \scriptsize
        \alertbf{Test}\par\smallskip
        \texttt{test-driven-development}
        \par\medskip
        Turn each acceptance line into a failing test, then build to green.
      \end{tdblock}
    \end{column}
    \begin{column}{0.23\textwidth}
      \begin{tdblock}
        \scriptsize
        \alertbf{Review}\par\smallskip
        \texttt{code-review}
        \par\medskip
        Read the diff as a stranger against standards and spec.
      \end{tdblock}
    \end{column}
    \begin{column}{0.23\textwidth}
      \begin{tdblock}
        \scriptsize
        \alertbf{Diagnose}\par\smallskip
        \texttt{diagnosing-bugs}
        \par\medskip
        Trace the symptom into the pure core, fix the cause.
      \end{tdblock}
    \end{column}
    \begin{column}{0.23\textwidth}
      \begin{tdblock}
        \scriptsize
        \alertbf{Deepen}\par\smallskip
        \texttt{improve-codebase-architecture}
        \par\medskip
        Pull one clean boundary behind the core, proposal only.
      \end{tdblock}
    \end{column}
  \end{columns}
\end{frame}
\fi

% === test-driven-development ===
\begin{frame}{\texttt{test-driven-development}}
  \slidesources{beck2002tdd}
  \begin{withmargin}
    \begin{tdmain}
      Each acceptance line in \texttt{issues/} becomes a test that fails before it
      passes. The skill drives the red, green, refactor loop one small change at a
      time, running \texttt{node --test} after each. \alertbf{Write the failing test
      before the code that satisfies it}, so each feature earns
      a real acceptance.
      \vskip0.5em
      \begin{tdblock}
        {\ttfamily\small
        > test-driven-development on 01-banana-on-screen\\
        RED: banana not shown, test fails\\
        GREEN: banana renders and runs, test passes}
      \end{tdblock}
    \end{tdmain}
    \begin{tdmargin}
      Fix the code, (never) the test, and watch the runner execute. 
      This concept can obviously be extended to integration and end-to-end tests as well.
    \end{tdmargin}
  \end{withmargin}
\end{frame}

\iffalse
% === code-review ===
\begin{frame}{\texttt{code-review}}
  \slidesources{martin2008clean}
  \begin{withmargin}
    \begin{tdmain}
      Read the diff as if a stranger wrote it: drop all prior context so the agent
      cannot defend its own code, then make two passes.
      \begin{itemize}\itemsep4pt
        \item \textbf{Standards:} check the diff against \texttt{AGENTS.md} and
          \texttt{GLOSSARY.md}.
        \item \textbf{Spec:} check it against the issue acceptance line or the
          \texttt{PRD}.
        \item Keep the core in \texttt{game.js} free of canvas and DOM.
        \item \alertbf{Report the problems, do not rewrite the code.}
      \end{itemize}
    \end{tdmain}
    \begin{column}{0.30\textwidth}
      \begin{tdblock}
        {\ttfamily\small
        > run code-review on the diff\\
        STANDARDS: game.js imports canvas, violates AGENTS.md\\
        SPEC: acceptance line met}
      \end{tdblock}
    \end{column}
  \end{withmargin}
\end{frame}
\fi


% === diagnosing-bugs ===
\begin{frame}{\texttt{diagnosing-bugs}}
  \slidesources{feathers2004legacy}
  \begin{withmargin}
    \begin{tdmain}
      Bugs will show up. Fixing it is mandatory, but then 
      hardening your codebase, by fortifying the fix with tests is best.
      \begin{itemize}\itemsep4pt
        \item Reproduce the bug with a test.
        \item Fix and document the bug.
        \item Confirm the fix by rerunning your test.
        \item \alertbf{Add test to suite.}
      \end{itemize}
    \end{tdmain}
    \begin{column}{0.30\textwidth}
      \begin{tdblock}
        {\ttfamily\small
        > diagnosing-bugs on ``some error message''}
      \end{tdblock}
    \end{column}
  \end{withmargin}
\end{frame}

\iffalse
% === improve-codebase-architecture ===
\begin{frame}{\texttt{improve-codebase-architecture}}
  \slidesources{ousterhout2018philosophy,fowler2018refactoring}
  \begin{withmargin}
    \begin{tdmain}
      Decay creeps in when concerns bleed together, like physics tangled into the
      render loop in \texttt{index.html}. The skill reads \texttt{GLOSSARY.md} and
      scans for mixed concerns.
      \begin{itemize}\itemsep4pt
        \item Pick \alertbf{one clean boundary} to pull into the pure core.
        \item Name what moves behind \texttt{game.js}, proposal only.
        \item Never a sweeping rewrite, never a direct edit.
        \item Keep \texttt{createGame}, \texttt{nextState}, \texttt{collides} a deep
          module you can test.
      \end{itemize}
    \end{tdmain}
    \begin{column}{0.30\textwidth}
      \begin{tdblock}
        {\ttfamily\small
        > run improve-codebase-architecture\\
        physics tangled into render loop in index.html\\
        proposal: move step() behind nextState in game.js}
      \end{tdblock}
    \end{column}
  \end{withmargin}
\end{frame}
\fi


% ======================================================================
% C. BREAKOUT -- vibe the banana again, now tested.
% ======================================================================

% === Vibe an infinitely running banana ===
\begin{frame}{\alertbf{Vibe} an infinitely running banana}
  \sideimagecover[0.33]{banana_1}
  \sidecaption{The banana runner, now building.}
  \begin{sidebody}
    Same banana, now under the discipline of tests. Take one acceptance line
    from the plan and run the whole build chain on it, back to back.
    \begin{itemize}
      \item \texttt{test-driven-development}. Write a test, and a feature against it.
      \item \texttt{diagnosing-bugs}: Encounter a bug, fix it, harden against it.
    \end{itemize}
    \vskip0.4em
    The deliverable is a \alertbf{tested banana core} you trust without reading
    every line. And what remains is taste.
  \end{sidebody}
\end{frame}


% References are collected in one shared package at the end of the deck
% (see disciplined_vibe.tex), not per unit.
